
\exo

Sans calcul, justifier les inégalités suivantes :

\vspace{-6pt}
\setlength{\columnsep}{1cm}
\setlength{\columnseprule}{0pt}
\begin{multicols}{2}
\begin{enumerate}
\item $0,5^2 < 0,75^2$.
\medskip
\item $\big(\frac{1}{3}\big)^2 > \big(\frac{1}{4}\big)^2$.
\item $(-3)^2 > (-1)^2$.
\medskip
\item $(-0,1)^2 < (-0,2)^2$.
\end{enumerate}
\end{multicols}

\exo

Soit $x$ un nombre réel. Que peut-on dire de $x^2$, dans chacun des cas suivants ?
\vspace{-6pt}
\setlength{\columnsep}{1cm}
\setlength{\columnseprule}{0pt}
\begin{multicols}{3}
\begin{enumerate}
\item $x\geqslant 5$.
\medskip
\item $x> 3$.
\medskip
\item $x\leqslant -2$.
\item $x< -7$.
\medskip
\item $1<x<100$.
\medskip
\item $-5<x<-2$.
\item $-10\leqslant x\leqslant 0$.
\medskip
\item $-5\leqslant x\leqslant 2$.
\medskip
\item $-5<x<5$.
\end{enumerate}
\end{multicols}

\exo

\begin{minipage}[t]{11.5cm}
Alice a remarqué que, sur le dessin ci-contre, les côtés des trois carrés mesurent tous entre $13,4$ et $13,6$ centimètres.

\begin{enumerate}
\item Déterminer un encadrement de l'aire d'un carré.
\item En déduire un encadrement de la somme des aires de ces trois carrés.
\item Alice souhaite agrandir son dessin, en multipliant toutes les longueurs par $10$. Elle sait qu'une bombe de peinture lui permet de couvrir environ \SI[retain-explicit-plus]{2}{\m^2}.

Combien de bombes de peinture lui faudra-t-il ?
\end{enumerate}
\end{minipage}
\hfill
\begin{minipage}[t]{5.5cm}
~

\vspace{-12pt}
\def\xmin{-1}
\def\xmax{2.5}
\def\ymin{-0.1}
\def\ymax{2.5}
\def\xunite{1.5cm}
\def\yunite{1.5cm}
\psset{unit=\xunite, xunit=\xunite, yunit=\yunite, algebraic=true, dimen=middle, dotstyle=*, dotsize=3pt 0, linewidth=0.8pt, arrowsize=3pt 2, arrowinset=0.25, comma=true}
\begin{pspicture*}(\xmin,\ymin)(\xmax,\ymax)
%\psgrid[xunit=\xunite, yunit=\yunite, subgriddiv=0, gridlabels=0, gridcolor=lightgray](0,0)(\xmin,\ymin)(\xmax,\ymax)
%\psaxes[labelFontSize=\scriptstyle, xAxis=true, yAxis=true, Dx=1, Dy=1, ticksize=-2pt 0, subticks=1]{->}(0,0)(\xmin,\ymin)(\xmax,\ymax)[{\footnotesize $x$},135][{\footnotesize $y$},-45]
\pspolygon[fillcolor=black,fillstyle=solid,opacity=0.1](0,0)(1,0)(1,1)(0,1)
\psrotate(1,1.2){-20}{\pspolygon[fillcolor=black,fillstyle=solid,opacity=0.1](1,1.2)(2,1.2)(2,2.2)(1,2.2)}
\psrotate(-0.1,1.1){30}{\pspolygon[fillcolor=black,fillstyle=solid,opacity=0.1](-0.1,1.1)(0.9,1.1)(0.9,2.1)(-0.1,2.1)}
\end{pspicture*}
\end{minipage}

\exo

\begin{minipage}[t]{11.5cm}
Un styliste imagine un foulard carré, dont le côté est compris entre $45$ et $50$ centimètres.

Il souhaite que les quatre carrés colorés aient la même aire, et que la surface colorée représente \SI[retain-explicit-plus]{80}{\percent} de la surface totale.

Déterminer un encadrement de la longueur des côtés des carrés.
\end{minipage}
\hfill
\begin{minipage}[t]{5.5cm}
~

\vspace{-12pt}
\def\xmin{-1}
\def\xmax{3}
\def\ymin{-0.1}
\def\ymax{2.5}
\def\xunite{1.2cm}
\def\yunite{1.2cm}
\psset{unit=\xunite, xunit=\xunite, yunit=\yunite, algebraic=true, dimen=middle, dotstyle=*, dotsize=3pt 0, linewidth=0.8pt, arrowsize=3pt 2, arrowinset=0.25, comma=true}
\begin{pspicture*}(\xmin,\ymin)(\xmax,\ymax)
%\psgrid[xunit=\xunite, yunit=\yunite, subgriddiv=0, gridlabels=0, gridcolor=lightgray](0,0)(\xmin,\ymin)(\xmax,\ymax)
%\psaxes[labelFontSize=\scriptstyle, xAxis=true, yAxis=true, Dx=1, Dy=1, ticksize=-2pt 0, subticks=1]{->}(0,0)(\xmin,\ymin)(\xmax,\ymax)[{\footnotesize $x$},135][{\footnotesize $y$},-45]
\pspolygon[fillcolor=black,fillstyle=solid,opacity=0.1](0,0)(1,0)(1,1)(0,1)
\pspolygon[fillcolor=black,fillstyle=solid,opacity=0.1](1.5,0)(2.5,0)(2.5,1)(1.5,1)
\pspolygon[fillcolor=black,fillstyle=solid,opacity=0.1](0,1.5)(1,1.5)(1,2.5)(0,2.5)
\pspolygon[fillcolor=black,fillstyle=solid,opacity=0.1](1.5,1.5)(2.5,1.5)(2.5,2.5)(1.5,2.5)
\pspolygon(0,0)(2.5,0)(2.5,2.5)(0,2.5)
\end{pspicture*}
\end{minipage}

\exo

Vrai ou Faux ? Justifier la réponse.

\begin{enumerate}
\item Les nombres $7$ et $-7$ ont la même image par la fonction carré.
\smallskip
\item Le nombre $\sqrt{3}$ est le seul antécédent de $3$ par la fonction carré.
\smallskip
\item Le nombre $-2$ est un antécédent de $-4$ par la fonction carré.
\end{enumerate}

\exo

\begin{minipage}[t]{11.5cm}
La courbe $\cC$ représente la fonction carré sur l'intervalle $\intff{-2}{2}$ et la droite $d$ est constituée des points d'ordonnée $1$.
\begin{enumerate}
\item
\begin{enumerate}
\item Donner les abscisses des points d'intersection de $\cC$ et $d$.
\item En déduire les solutions de l'équation $x^2=1$ dans l'intervalle $\intff{-2}{2}$ ?
\end{enumerate}
\item
\begin{enumerate}
\item Repasser en couleur la partie de $\cC$ située \og en-dessous \fg{} de la droite $d$.
\item En déduire les solutions dans $\intff{-2}{2}$ de l'inéquation $x^2<1$.
\end{enumerate}
\end{enumerate}
\end{minipage}
\hfill
\begin{minipage}[t]{5.5cm}
~

\vspace{-12pt}
\def\xmin{-3}
\def\xmax{3}
\def\ymin{-1}
\def\ymax{5}
\def\xunite{0.8cm}
\def\yunite{0.7cm}
\psset{unit=\xunite, xunit=\xunite, yunit=\yunite, algebraic=true, dimen=middle, dotstyle=*, dotsize=3pt 0, linewidth=0.8pt, arrowsize=3pt 2, arrowinset=0.25, comma=true}
\begin{pspicture*}(\xmin,\ymin)(\xmax,\ymax)
\psgrid[xunit=\xunite, yunit=\yunite, subgriddiv=0, gridlabels=0, gridcolor=lightgray](0,0)(\xmin,\ymin)(\xmax,\ymax)
\psaxes[labelFontSize=\scriptstyle, xAxis=true, yAxis=true, Dx=1, Dy=1, ticksize=-2pt 0, subticks=1]{->}(0,0)(-2.99,-0.99)(\xmax,\ymax)[{\footnotesize $x$},135][{\footnotesize $y$},-45]
\psplot[linewidth=1.2pt, plotpoints=200]{-2}{2}{x*x}
\psplot[linewidth=1.2pt, plotpoints=200]{\xmin}{\xmax}{1}
\begin{footnotesize}
\uput{4pt}[000](2,4){$\cC$}
\uput{4pt}[090](2.5,1){$d$}
\end{footnotesize}
\end{pspicture*}
\end{minipage}

\exo

On note $a$ et $b$ deux réels, et on considère la proposition : \og Si $a<b$, alors $a^2<b^2$ \fg{}.
\begin{enumerate}
\item Cette proposition est-elle vraie ou fausse ? Justifier.
\item Énoncer la proposition réciproque. Est-elle vraie ou fausse ? Justifier.
\end{enumerate}

\exo

Reprendre l'exercice précédent avec la proposition : \og Si $x^2=4$, alors $x=2$ \fg{}.

\exo

Résoudre dans $\R$ les équations et inéquations :

\vspace{-6pt}
\setlength{\columnsep}{1cm}
\setlength{\columnseprule}{0pt}
\begin{multicols}{4}
\begin{enumerate}
\item $x^2=5$.
\medskip
\item $x^2<5$.
\medskip
\item $x^2\geqslant 5$.
\item $x^2=2$.
\medskip
\item $x^2\leqslant 2$.
\medskip
\item $x^2>2$.
\item $x^2=-3$.
\medskip
\item $x^2<-3$.
\medskip
\item $x^2\geqslant -3$.
\item $x^2\leqslant 8$.
\medskip
\item $x^2\geqslant 4$.
\medskip
\item $x^2>2$.
\end{enumerate}
\end{multicols}
\vspace{-6pt}

\exo

Construire le tableau de signe des expressions suivantes :

\vspace{-6pt}
\setlength{\columnsep}{1cm}
\setlength{\columnseprule}{0pt}
\begin{multicols}{3}
\begin{enumerate}
\item $P(x)=(x-1)(x-2)$.
\medskip
\item $P(x)=(x-4)(2-x)$.
\medskip
\item $P(x)=(-x+5)(2x+1)$.
\medskip
\item $P(x)=(7-x)(10-x)$.
\medskip
\item $P(x)=x(x-1)(2x-1)$.
\medskip
\item $P(x)=x^2(x+3)$.
\medskip
\item $P(x)=(x^2+1)(5-10x)$.
\medskip
\item $P(x)=(-2x+6)(4x-8)$.
\medskip
\item $P(x)=(x^2-7)(3x+9)$.
\medskip
\end{enumerate}
\end{multicols}
\vspace{-6pt}

\exo

Résoudre dans $\R$ les inéquations suivantes :

\vspace{-6pt}
\setlength{\columnsep}{1cm}
\setlength{\columnseprule}{0pt}
\begin{multicols}{3}
\begin{enumerate}
\item $(x-1)(x-2)>0$.
\medskip
\item $(x-4)(2-x)\leqslant 0$.
\medskip
\item $(-x+5)(2x+1)>0$.
\item $(7-x)(10-x)<0$.
\medskip
\item $x(x-1)(2x-1)\geqslant 0$.
\medskip
\item $x^2(x+3)\leqslant 0$.
\item $(x^2+1)(5-10x)>0$.
\medskip
\item $(-2x+6)(4x-8)>0$.
\medskip
\item $(x^2-7)(3x+9)\leqslant 0$.
\end{enumerate}
\end{multicols}
\vspace{-6pt}

\exo

Résoudre dans $\R$ les équations et inéquations suivantes :

\vspace{-6pt}
\setlength{\columnsep}{1cm}
\setlength{\columnseprule}{0pt}
\begin{multicols}{4}
\begin{enumerate}
\item $5x^2\geqslant 2x$.
\item $-8x^2> 2x$.
\item $-2x^2\geqslant 10x$.
\item $7x^3< 3x$.
\end{enumerate}
\end{multicols}
\vspace{-6pt}

\exo

On étudie la fonction $d$ qui, à la vitesse $v$ d'un véhicule (exprimée en \si[per-mode=symbol]{\m\per\s}) associe la distance de freinage $d(v)$ (exprimée en \si{\m}). Cette fonction est définie sur $\intfo{0}{+\infty}$ par $d(v)=k\times v^2$, où $k$ est un coefficient qui dépend notamment de l'état de la route.

\begin{enumerate}
\item Sur route sèche, la valeur de $k$ est égale à $0,08$.
\begin{enumerate}
\item Montrer qu'une vitesse de \SI[per-mode=symbol]{80}{\km\per\hour} est environ égale à \SI[per-mode=symbol]{22,22}{\m\per\s}.

En déduire la distance de freinage sur route sèche pour une vitesse égale à \SI[per-mode=symbol]{80}{\km\per\hour}.
\item À partir de quelle vitesse (en \si[per-mode=symbol]{\km\per\hour}) la distance de freinage sur route sèche dépasse-t-elle \SI{45}{\m} ?
\end{enumerate}
\end{enumerate}
\begin{minipage}[t]{10.5cm}
\begin{enumerate}[start=2]
\item Sur route mouillée, la valeur de $k$ est différente de $0,08$.

La représentation graphique ci-contre donne la distance de freinage sur route mouillée en fonction de la vitesse.
\begin{enumerate}
\item En utilisant cette représentation, déterminer la vitesse (en \si[per-mode=symbol]{\km\per\hour}) pour laquelle la distance de freinage sur route mouillée est égale \SI{100}{\m}.
\item Déterminer la valeur de $k$ sur route mouillée.
\end{enumerate}
\end{enumerate}
\end{minipage}
\hfill
\begin{minipage}[t]{6.5cm}
~

%\vspace{-6pt}
\def\xmin{-5}
\def\xmax{35}
\def\ymin{-15}
\def\ymax{160}
\def\xunite{0.15cm}
\def\yunite{0.035cm}
\psset{unit=\xunite, xunit=\xunite, yunit=\yunite, algebraic=true, dimen=middle, dotstyle=*, dotsize=4pt 0, linewidth=0.8pt, arrowsize=3pt 2, arrowinset=0.25, comma=true}
\begin{pspicture*}(\xmin,\ymin)(\xmax,\ymax)
\multido{\n=0+5}{8}{\psline[linecolor=lightgray](\n,0)(\n,\ymax)}
\multido{\n=0+20}{9}{\psline[linecolor=lightgray](0,\n)(\xmax,\n)}
%\psgrid[xunit=\xunite, yunit=\yunite, Dx=5, Dy=20, subgriddiv=0, gridlabels=0, gridcolor=lightgray](0,0)(\xmin,\ymin)(\xmax,\ymax)
\psaxes[labelFontSize=\scriptstyle, xAxis=true, yAxis=true, Dx=5, Dy=20, ticksize=-2pt 0, subticks=1]{->}(0,0)(0,0)(\xmax,\ymax)[{\footnotesize Vitesse (en \si[per-mode=symbol]{\m\per\s})},135][{\footnotesize Distance (en \si{\m})},-45]
\psplot[linewidth=1.2pt, plotpoints=200]{0}{30}{0.16*x^2}
\psplot[plotstyle=dots,plotpoints=7]{0}{30}{0.16*x^2}
%\begin{footnotesize}
%\uput{3pt}[315](1.4,2){$\cP$}
%\uput{3pt}[90](2.7,5){$\Delta$}
%\uput{3pt}[270](1,0){$1$}
%\uput{3pt}[180](0,1){$1$}
%\end{footnotesize}
\end{pspicture*}
\end{minipage}

